\documentclass[output=paper]{langscibook}
\ChapterDOI{10.5281/zenodo.22880278}
\title{\textit{A goroto} – The turtle}
\author{Naphtali Maion and Ondria Tavagaga}
\abstract{}
\begin{document}
\maketitle\label{chap:26}
\begin{figure}[H]
    \captionsetup[subfigure]{oneside,margin={0cm,0cm}}
    \subfloat[]{\includegraphics[trim=0 0 0 0,clip,width=.4\linewidth]{figures/26_1.png}\label{fig:26_1}}\hfil
    \subfloat[]{\includegraphics[trim=0 0 0 0,clip,width=.6\linewidth]{figures/26_2.png}\label{fig:26_2}}
    \caption{(a) \textit{Unava} -- A small kind of turtle, (b) \textit{Toorue} -- The largest kind of Turtle}
\end{figure}
\begin{xltabular}{\textwidth}{lXX}
	(\ref{ex:text19-1})&\textit{\textbf{Goroto}}&\textbf{Turtle}\\
	(\ref{ex:text19-2})&\textit{A goroto a iana a beera sana.}&The turtle is a very big swimming animal.\\
	(\ref{ex:text19-3})&\textit{Bean tara e, ean repaa tara voen:}&When you see it, you see:\\
	(\ref{ex:text19-4})&\textit{A goroto na tei me nana bono kahoo, o dorana o pakapaka,}&The turtle has a head and four flippers.\\
	(\ref{ex:text19-5})&\textit{O kapa vaa toonunae o kikis, o kapa vaa domonae saka kikis haana.}&The shell on its back is hard, the shell on its chest is not hard.\\
	(\ref{ex:text19-6})&\textit{Teo dorana o pakapaka, o buaku vaa kahoo nae na dao raara bono huahua.}&As for the four flippers, the two flippers near its head we call the paddles.\\
	(\ref{ex:text19-7})&\textit{A goroto na huahua vakikis ni nana bono buo pakapaka bona.}&The turtle strongly paddles with these two flippers.\\
	(\ref{ex:text19-8})&\textit{O buo pakapaka vaa murinae to vakarovo bata ni nae bara na hua ni nae mee.}&The two hind flippers it uses for steering and also for paddling.\\
	(\ref{ex:text19-9})&\textit{A goroto na vaavuriha nana mee be tavusu kahi bona tahii, ore paa ruuvu gunaha vahaa.}&The turtle also breathes when it comes out of the water, then it dives down again.\\
	(\ref{ex:text19-10})&\textit{A goroto vai a moon be tau vahuhu, ore paa gunaha vo teo kakaamuru.}&When the female turtle is about to lay eggs, it goes ashore to the white sand.\\
	(\ref{ex:text19-11})&\textit{Na kasi nana bona govee kasuana}&It digs a hole in the ground,\\
	(\ref{ex:text19-12})&\textit{ore paa vahuhu komana govee}&then it lays its eggs into the hole\\
	(\ref{ex:text19-13})&\textit{ore paa pae vahaa bona.}&and fills the hole up again.\\
	(\ref{ex:text19-14})&\textit{Evehee bea goroto pae vahaa bona govee, na varovarova nae,}&But when the turtle fills up the hole again, it becomes invisable,\\
	(\ref{ex:text19-15})&\textit{a aba ore saka antee haa bona tea vurahe vakahu.}&so that a human being cannot quickly find it.\\
	(\ref{ex:text19-16})&\textit{A rupi a bero kurusu to vahuhu nana a goroto.}&The turtle lays many eggs indeed.\\
	(\ref{ex:text19-17})&\textit{Toro oha nana bono taonim saavun.}&It must be more than fifty.\\
	(\ref{ex:text19-18})&\textit{Bea peha aba naonao tea ruuvu goroto, na nao vatotobin vo nana irihi.}&When someone goes diving for turtles, he goes straight to the edge of the reef.\\
	(\ref{ex:text19-19})&\textit{O goroto na teitei rori irihi tea maa avaava vasu ei teo hum namana.}&The turtles stay at the edge of the reef in the caves of the rocks at the side of the ocean.\\
	(\ref{ex:text19-20})&\textit{Bea tahii masi, ore paa vurahe vakahu koa  bono goroto}&When the sea is at low tide (so that the reef is visable), then he can quickly find the turtles,\\
	(\ref{ex:text19-21})&\textit{beori ani rori bono vasukare bara tea maa avaava vasu ae tea maa vaagoroto.}&when they are eating the \textit{vasukare} corals and are in the rock caves and in the turtle tunnels.\\
\end{xltabular}
\begin{xltabular}{\textwidth}{lXX}
	(\ref{ex:text19-22})&\textit{\textbf{Havee to gono kaakoo mi riori bona goroto}}&\textbf{How they catch a turtle}\\
	(\ref{ex:text19-23})&\textit{A goroto na ato riori tea kave,}&They get hold of the turtle with the net,\\
	(\ref{ex:text19-24})&\textit{ei beori nao tea veo goroto.}&this means when they go turtle netting.\\
	(\ref{ex:text19-25})&\textit{A kave bona a beera sana.}&This net is very big.\\
	(\ref{ex:text19-26})&\textit{Na antee nana tea kana bono buaku saavun goroto.}&It can catch twenty turtles.\\
	(\ref{ex:text19-27})&\textit{Amaa peha maa vuri to naonao riori tea ruuvu goroto.}&Sometimes they go diving for turtles.\\
	(\ref{ex:text19-28})&\textit{A goroto mee na kapa riori bea goroto gunaha nao tea vahuhu kasuana.}&They also catch a turtle when the turtle goes ashore to lay its eggs in the sand.\\
	(\ref{ex:text19-29})&\textit{Tea maa vuri sunuu, a goroto mee na vurahe riori ei be masi pina ni tea maa kopua.}&When they go fishing with torches, they also find the turtle when it gets trapped in the lagoons at low tide.\\
\end{xltabular}
\begin{xltabular}{\textwidth}{lXX}
	(\ref{ex:text19-30})&\textit{\textbf{Havee to kaku goroto vaavaha vo ria abana vaa Teapu}}&\textbf{How turtles are butchered by the Teop people}\\
	(\ref{ex:text19-31})&\textit{Nabunuu roho, a kaku goroto, amaa abana bebeera koa to pakupaku roho bari.}&In former times the butchering of turtles was only done by the big men.\\
	(\ref{ex:text19-32})&\textit{Ean saka antee haa nom tea tara ta peha beiko ge ta peha visoasi teo hum vai to kaku goroto riori.}&You could not see any child or any young man at the place where they were butchering turtles.\\
	(\ref{ex:text19-33})&\textit{A kiu tea abana bebeera koa, paku vonae tea kaku guu mee.}&It was exclusively the work of the chiefs, it was like butchering pigs.\\
	(\ref{ex:text19-34})&\textit{Nomana amaa matapaku vai paa rova kahi avuaraara.}&Today these practices have disappeared from us.\\
	(\ref{ex:text19-35})&\textit{Nomana bean tau kaku a goroto -- paku koai; ahiki vira to moroko nae vahaa.}&Today, when you are about to butcher a turtle -- just do it; indeed, there is not any talk about it.\\
\end{xltabular}
\begin{xltabular}{\textwidth}{lXX}
	(\ref{ex:text19-36})&\textit{\textbf{Amaa hanana vaa tea kaku goroto:}}&\textbf{The ways of butchering turtles:}\\
	(\ref{ex:text19-37})&\textit{Rahi gunaha a goroto teo hum to kahi kaku noman e.}&Pull the turtle ashore to the place where you are going to butcher it.\\
	(\ref{ex:text19-38})&\textit{Garuhu vira nio kasuana.}&Thoroughly wash the sand off.\\
	(\ref{ex:text19-39})&\textit{Navuhu vamate a goroto, navuhu kii gogoonae.}&Hit the turtle dead, hit it on its throat.\\
	(\ref{ex:text19-40})&\textit{Pee o kapa vaa domonae, pee murina koa o vavaarihi.}&Cut the shell on its chest, subsequently cut the sides.\\
	(\ref{ex:text19-41})&\textit{Ean repaa kapae.}&Then you take its shell off.\\
	(\ref{ex:text19-42})&\textit{Rake nao o pusu no hiapa to taneo maana teo orooromo,}&Search for the upper end of the intestines that begin in its throat,\\
	(\ref{ex:text19-43})&\textit{pee kurukurusui, rahi tavusu kahi bona komana goroto.}&cut them off, pull them out of the belly of the turtle.\\
	(\ref{ex:text19-44})&\textit{Rahi tavusu vakavara o hiapa ae amaa meha taba vaa tea goroto,}&When you have completely pulled out the instestines and the other things in the turtle,\\
	(\ref{ex:text19-45})&\textit{kana tavusu o revasin. }&scoop the blood out.\\
	(\ref{ex:text19-46})&\textit{Peepee vihu amaa pakapaka nae (o dorana) bara o kahoo.}&Cut off its flippers (the four) and the head.\\
	(\ref{ex:text19-47})&\textit{Vaanoto o suraa o beera mene rava bona mata goroto.}&Light a big fire for heating the pieces of the turtle.\\
	(\ref{ex:text19-48})&\textit{Rava vakavara, vaatei a mata goroto teo bana.}&After having heated them up, put the turtle pieces on a mat.\\
	(\ref{ex:text19-49})&\textit{Vaagunagunaha a mata goroto, ean repaa peepee vaantee a peepeha mono goroto.}&Cut the turtle pieces further up, then you cut enough for each turtle parcel.\\
	(\ref{ex:text19-50})&\textit{Monomono a goroto. Kavara vai voen, vaatei vai tea tao.}&Wrap the turtle up into parcels. When it's finished in this way, put (the parcels) into the earth oven.\\
	(\ref{ex:text19-51})&\textit{Amaa hum vaa komana goroto na aniani ria ta Teapu.}&The parts inside the turtle are eaten by the Teop people\\
	(\ref{ex:text19-52})&\textit{O hiapa na peepee vakakato riori}&The intestines are cut into short pieces,\\
	(\ref{ex:text19-53})&\textit{ore paa siri kopaa bari,}&then they are torn into strips,\\
	(\ref{ex:text19-54})&\textit{ore paa garuhu vamataa bari,}&then properly washed\\
	(\ref{ex:text19-55})&\textit{ore paa vaatei bari tea nahu.}&and put into a pot.\\
	(\ref{ex:text19-56})&\textit{O gonogonono mee na peepee riori, ore paa vaatei bari tea nahu.}&The lungs are also cut up and put into a pot.\\
	(\ref{ex:text19-57})&\textit{Amaa hum vai mee, a koburu, o ganusu ae amaa meha hum.}&Also the following parts, the heart, cartilages, and the other parts\\
	(\ref{ex:text19-58})&\textit{Be kavara, ore paa vaahio bona nahu teo suraa.}&When it is finished, the pot is set on the fire.\\
	(\ref{ex:text19-59})&\textit{Na nahunahu mi riori bono hiihii ore nin mataa.}&It is cooked with ginger so that it tastes good.\\
\end{xltabular}
\end{document}